\chapter{Cross-Cutting Concerns}
\label{ch:crosscutting}

\section{Idempotency}

Idempotency is enforced at two levels:

\begin{enumerate}
    \item \textbf{Outbox side:} each \texttt{OutboxMessage} carries a unique \texttt{EventId} (a \texttt{Guid} generated at write time). Consumers use this field as their deduplication key.
    \item \textbf{Consumer side:} \texttt{InboxDeduplicationService} checks \texttt{ProcessedInboxMessage} before any handler executes. If the \texttt{EventId} has already been recorded, the message is discarded. This makes all handlers idempotent by construction, regardless of their internal logic.
\end{enumerate}

\section{Error Handling and Dead-Letter Strategy}

\begin{table}[H]
\centering
\caption{Error handling strategy by failure point}
\label{tab:error-handling}
\begin{tabularx}{\textwidth}{@{} L{3cm} L{4cm} X @{}}
\toprule
\textbf{Failure point} & \textbf{Behaviour} & \textbf{Operator action} \\
\midrule
Outbox write fails & Order placement rolls back (same transaction) & None — customer sees failure, can retry \\
\addlinespace
RabbitMQ publish fails & \texttt{Attempts} incremented; retried on next dispatcher tick & Monitor \texttt{Attempts} column; investigate if broker is down \\
\addlinespace
\texttt{Attempts > N} & Row stays in outbox as dead-letter record & Manual intervention: inspect payload, republish or void \\
\addlinespace
Consumer handler throws & Message NACK'd to RabbitMQ; requeued or routed to DLQ after configured retries & Inspect RabbitMQ DLQ; fix handler bug; republish \\
\addlinespace
Duplicate message received & Silently discarded by \texttt{Inbox\-Deduplication\-Service} & None required \\
\bottomrule
\end{tabularx}
\end{table}

\section{Observability}

The system provides three observable surfaces without requiring external tooling:

\begin{description}
    \item[Outbox table] Every cross-context message has a row with \texttt{OccurredOnUtc}, \texttt{DispatchedOnUtc}, and \texttt{Attempts}. The lag between these two timestamps is the propagation delay. Rows with \texttt{DispatchedOnUtc = NULL} and high \texttt{Attempts} signal delivery problems.
    \item[ProcessedInboxMessage table] An audit log of every inbound message successfully processed, with timestamps.
    \item[Order status transitions] The \texttt{OrderStatus} field records the current state. Combined with the nopCommerce order notes (added by each handler), this constitutes a state transition audit trail.
    \item[RabbitMQ management dashboard] Exchange throughput, queue depths, and DLQ contents are visible in the RabbitMQ management UI (\texttt{http://localhost:15672} in the Docker Compose environment).
\end{description}

\section{Security}

\begin{itemize}
    \item RabbitMQ credentials are passed via environment variables (not hardcoded). In the Docker Compose environment they are set in \texttt{.env}.
    \item The outbox table is inside the nopCommerce application database, protected by the same SQL Server access controls as all other tables.
    \item No new HTTP endpoints are exposed by the two plugins — all communication is DB reads/writes and AMQP.
    \item The event payload is JSON. No sensitive customer data (payment card numbers, passwords) is included in any cross-context event. Payment operations use nopCommerce's existing payment service abstractions internally.
\end{itemize}

\section{Scalability Constraints}

The current design is explicitly single-instance. The in-memory nopCommerce cache (\texttt{IStaticCacheManager}) is not shared across instances. Running two nopCommerce processes would cause cache divergence and double-dispatching from the outbox. Horizontal scaling would require Redis (or another distributed cache) and a leader-election mechanism for the dispatcher task. This is acknowledged as out of scope for the demo.

